\customchapter{CONCLUSIONES}

Esta fase del trabajo implementó y evaluó las unidades U1 a U3 de la
arquitectura propuesta: la construcción del entorno de simulación, un
agente de control semafórico basado en Q-learning tabular y la
instrumentación de las métricas de desempeño, sobre cuatro escenarios
sintéticos de complejidad creciente y con un protocolo de evaluación
que compara tres controles sobre las mismas diez semillas. De los
resultados del Capítulo~\ref{cap:resultados} se desprenden las
siguientes conclusiones.

\noindent\textbf{1. El agente tabular no alcanza las metas planteadas y
solo mejora al control de tiempo fijo en el escenario más simple.} Con
la política congelada y diez semillas de evaluación, el retraso
promedio por vehículo baja un 6.4\,\% en el cruce de dos fases y sube
entre un 6 y un 17\,\% en los otros tres escenarios. La meta declarada
de una reducción del 25 al 40\,\% no se cumple en ningún caso. La
reducción de colas del 20\,\% se alcanza únicamente en el cruce simple
y en la media, con un intervalo de confianza que no excluye ese umbral.

\noindent\textbf{2. La función de recompensa no está alineada con la
métrica que la tesis declara como primaria.} La recompensa penaliza
vehículos detenidos y su tiempo de espera, y el agente la mejora de
forma consistente; pero la meta se formula sobre el retraso, que
incluye además frenadas y aceleraciones. La consecuencia es medible: en
el cruce simple la espera detenida baja un 26\,\% mientras el retraso
solo baja un 6\,\%, y en el corredor la espera baja un 13\,\% mientras
el retraso sube un 6\,\%. El agente aprende a cortar el verde con
frecuencia, hasta un 77\,\% más de cambios de fase que el control fijo,
lo que fragmenta las colas pero aumenta las paradas por vehículo en los
cuatro escenarios. Optimizar una señal correlacionada con el objetivo
no equivale a optimizar el objetivo.

\noindent\textbf{3. La representación tabular tiene un techo
identificable al crecer el número de fases.} En el cruce de dos fases el
agente visita 130 de 512 estados posibles y converge; en el cruce con
giros protegidos, donde el espacio se duplica a 1024 estados, visita
390 y su política congelada resulta peor que el control fijo, con una
dispersión entre semillas casi veinte veces mayor. Aumentar el número
de episodios no corrige el problema, porque la causa no es falta de
entrenamiento sino la fracción del espacio de estados que queda sin
visitar. Este resultado es el argumento empírico para pasar a métodos
con aproximación de funciones, como DQN y PPO.

\noindent\textbf{4. Los agentes sin comunicación no reproducen el
beneficio de la coordinación.} En la red de tres intersecciones, tres
agentes independientes quedan un 13\,\% por detrás de un programa fijo
con onda verde, aunque superan ampliamente al mismo programa sin
desfases, que da 41.6 s de retraso frente a 25.1 s del agente. La
diferencia no está en la capacidad de reaccionar a la cola local, sino
en la anticipación que aporta coordinar las fases de intersecciones
sucesivas. Es la evidencia empírica que justifica el módulo U4.

\noindent\textbf{5. El control actuado es una referencia más exigente
que el control de tiempo fijo y la literatura revisada lo omite.} Con
las mismas fases y los mismos límites de verde que el agente, el
control actuado nativo del simulador reduce el retraso entre un 4 y un
34\,\% frente al tiempo fijo y supera al agente tabular en los cuatro
escenarios. Ninguno de los tres trabajos analizados en el estado del
arte lo emplea como línea base. Cualquier
resultado de control semafórico por aprendizaje por refuerzo que aspire
a ser relevante en la práctica debe medirse también contra él.

\noindent\textbf{6. El protocolo experimental resultó determinante para
la validez de las conclusiones.} Separar las semillas de entrenamiento
de las de evaluación, congelar la política antes de medir, usar las
mismas semillas en todos los controles y reportar intervalos de
confianza con pruebas pareadas cambió el signo de varias conclusiones
respecto de lo que sugerían las corridas con una sola semilla. Un caso
lo ilustra con claridad: en el corredor, dos entrenamientos que solo
diferían en el instante en que el agente toma sus decisiones dieron
resultados de signo opuesto, lo que muestra que una única corrida de
entrenamiento no caracteriza a la política aprendida.
